\section{Scenario \#3: ''Winchester'': May 20-27, 1862
}


\begin{scenariodata}
\textbf{Game Length:} 8 Turns\\
\end{scenariodata}

\begin{scenariodata}
(Weather: May 16-31 Column)\\
\end{scenariodata}

\begin{scenariodata}
\textbf{First Phase:} Confederates\\
\end{scenariodata}

\begin{scenariodata}
\textbf{Union:}\\
\end{scenariodata}

\begin{scenariodata}
Resource Factor 2 10\\
\end{scenariodata}

\begin{scenariodata}
\textbf{Initial Placement:}\\
\end{scenariodata}

\begin{scenariodata}
Hex A-16 (Romney) --: [51, 2A, 1C], (1F, 1S)\\
Hex B-25 (Moorefield) : [31, 1A], (1F, 1S)\\
Hex 8-37 (Franklin) : [31, 1C], 251, 2A, 2C,\\
\end{scenariodata}

\begin{scenariodata}
\textbf{and adjacent :} (Fremont, 1F, 2S, 1W)\\
Hex M-26 (Strasburg) : 131, 3A, 8C,\\
\textbf{and adjacent :} (Banks, 1S, 1W)\\
\end{scenariodata}

\begin{scenariodata}
Hex N-20 (Winchester) : 21, 2C, (1F, 4S, 2W)\\
Hex S-28 (Front Royal):\\
\textbf{and adjacent :} 31, (28, 1W)\\
\end{scenariodata}

\begin{scenariodata}
B\&O : [171, 3A, 3C], (1S)\\
\end{scenariodata}

\begin{scenariodata}
\textbf{Withdrawal Hexes:}\\
\end{scenariodata}

\begin{scenariodata}
N-7, S-10, Z-13, CC-21 through CC-33\\
(Supply/Wagon Entry Hexes:\\
\end{scenariodata}

\begin{scenariodata}
N-7, 5-10, Z-13, A-16, A-11, A-13,\\
CC-21 through CC-33)\\
\end{scenariodata}

\begin{scenariodata}
No additional Fortification counters.\\
\end{scenariodata}

\begin{scenariodata}
\textbf{Confederates:}\\
\end{scenariodata}

\begin{scenariodata}
\textbf{Resource Factor :} 10\\
\end{scenariodata}

\begin{scenariodata}
\textbf{Initial Placement:}\\
\end{scenariodata}

\begin{scenariodata}
Hex 0-37 (New Market) : 231,6A,\\
\end{scenariodata}

\begin{scenariodata}
\textbf{and adjacent :} (Jackson, 3S, 3W)\\
Hex T-43 (Conrad's Store) : 241, 4A,\\
\end{scenariodata}

\begin{scenariodata}
\textbf{and adjacent :} (Ewell, 25, 3W)\\
Within 9 hexes of -\\
\end{scenariodata}

\begin{scenariodata}
Hex 0-42 (Harrisonburg) : 9C, 1H, (Ashby)\\
Hex 0-51 (Staunton) > Tl, (1F, 1W)\\
\end{scenariodata}

\begin{scenariodata}
Hex 88-53 (Charlottesville) : 11, 1A, (1F)\\
(Partisans: McNeill, Gilmor)\\
\end{scenariodata}

\begin{scenariodata}
\textbf{Withdrawal Hexes:}\\
\end{scenariodata}

\begin{scenariodata}
CC-47 through CC-58\\
(Supply/Wagon Entry Hexes:\\
CC-47 through CC-58, 0-51, M-50)\\
\end{scenariodata}

\begin{scenariodata}
No additional Fortification counters.\\
\end{scenariodata}

\begin{scenariodata}
\textbf{Victory Conditions:} All other results are a draw.\\
\end{scenariodata}

\begin{scenariodata}
\textbf{ADVANCED GAME:} Union wins if they have more total Victory Points.\\
Confederates win if they have 50 or more fead in\\
\end{scenariodata}

total Victory Points.

\begin{scenariodata}
\textbf{BASIC GAME:} Same as the Advanced Game, only "spot'' the Con-\\
\end{scenariodata}

federates 10 Victory Points at the start of the game.

\begin{scenariodata}
\textbf{HISTORICAL RESULTS:} Confederate victory.\\
\end{scenariodata}

\begin{scenariodata}
\textbf{PLAYTESTER'S COMMENTS:} The Confederate strategic position is\\
\end{scenariodata}

very good -- a central position between the two separated masses of

\begin{scenariodata}
Union field armies. Either Union army can be attacked, preferably Banks,\\
\end{scenariodata}

as more Victory Points are available at the Northern end of the Mapboard. The Union player should try to move Fremont to a position where he can draw more supplies, then move to support Banks. Banks

should be careful not to get cut off, but try to withdraw as gradually as possible towards Harper's Ferry.

| een)

The Confederates moved swiftly, leaving only a cavalry screen to watch Fremont, and concentrating to smash the small Federa! force at Front Royal. This maneuver outflanked Banks at Strasburg, and forced him to fall back to defend the Potomac. The Confederates mauled the retreating Union columns near Winchester, although rain hindered the Confederate pursuit. Banks escaped with most of his army, although the quantity of supplies abandoned caused him to be known as "'Commis- sary' Banks to the Confederates. Jackson's maneuvers threatened Harper's Ferry, and the garrison there fell back across the Potomac to escape destruction. The victorious Confederates tore up the B\&O, and lingered near the Potomac, to appear to be a viable threat to Washington, D.C.

The news of the advance of Jackson's ''foot cavalry" electrified the

\begin{scenariodata}
South, coming as it did at a time when all other news was bad, and the\\
Union Army of the Potomac sat on the outskirts of Richmond. It caused\\
\end{scenariodata}

consternation in the North, and led to the dispatch of Irvin McDowell's Corps, which had been intended as the northern wing on the final advance on Richmond, to the Shenandoah. With Fremont to the west, the head of McDowell's Corps (Shield's division) to the east, and Banks to the north, Jackson found himself in a very precarious position...

